% =============================================================================
% 5.4 Pruebas Unitarias del Backend
% Archivo: 5.4_pruebas_unitarias_backend.tex
% =============================================================================

\section{Pruebas Unitarias del Backend}
\label{sec:pruebas-unitarias-backend}

Las pruebas unitarias representan el primer nivel de la pir\'amide de pruebas
de software y constituyen el mecanismo m\'as granular para verificar el
comportamiento correcto de las unidades individuales de c\'odigo. En el contexto
del servidor \textit{backend} de la plataforma, se implement\'o una estrategia
exhaustiva de pruebas unitarias que abarca las capas de enrutamiento, servicios
de negocio, componentes de seguridad, \textit{middleware} transversal y
funciones utilitarias.

\subsection{Resultados Generales de Ejecuci\'on}
\label{subsec:resultados-unitarias}

La ejecuci\'on completa de las pruebas unitarias del \textit{backend} arroj\'o
resultados plenamente satisfactorios, con la totalidad de los suites y casos
de prueba aprobados sin excepciones. La Tabla~\ref{tab:resumen-unitarias}
presenta el resumen cuantitativo de la ejecuci\'on.

\begin{table}[htbp]
\centering
\caption{Resumen de ejecuci\'on de pruebas unitarias del \textit{backend}}
\label{tab:resumen-unitarias}
\begin{tabular}{|l|c|}
\hline
\textbf{Par\'ametro} & \textbf{Valor} \\
\hline
\textit{Framework} de pruebas         & Jest (TypeScript) \\
\hline
Suites de pruebas ejecutados          & 16 \\
\hline
Suites aprobados                      & 16 (100\%) \\
\hline
Suites fallidos                       & 0 \\
\hline
Casos de prueba ejecutados            & 334 \\
\hline
Casos de prueba aprobados             & 334 (100\%) \\
\hline
Casos de prueba fallidos              & 0 \\
\hline
Tiempo total de ejecuci\'on           & 19.375 s \\
\hline
\end{tabular}
\end{table}

\subsection{Distribuci\'on de Suites por Categor\'ia}
\label{subsec:distribucion-suites}

Los diecis\'eis suites de pruebas se organizan en cinco categor\'ias funcionales
que reflejan la arquitectura en capas del servidor \textit{backend}. Esta
clasificaci\'on permite identificar con precisi\'on la cobertura de cada
componente arquitect\'onico y priorizar las \'areas que requieren
fortalecimiento en iteraciones futuras. La Tabla~\ref{tab:suites-categoria}
presenta la distribuci\'on detallada.

\begin{table}[htbp]
\centering
\caption{Distribuci\'on de suites de pruebas unitarias por categor\'ia}
\label{tab:suites-categoria}
\begin{tabular}{|l|p{6.5cm}|p{4.5cm}|}
\hline
\textbf{Categor\'ia} & \textbf{Suites incluidos} & \textbf{Alcance funcional} \\
\hline
Rutas (\textit{Routes}) & \texttt{iam.routes}, \texttt{sedes.routes}, \texttt{audit.routes}, \texttt{auth.routes}, \texttt{onboarding.routes} & Validaci\'on de endpoints HTTP, par\'ametros de entrada y c\'odigos de respuesta. \\
\hline
Servicios (\textit{Services}) & \texttt{email.service} & L\'ogica de negocio del servicio de correo electr\'onico y env\'io de notificaciones. \\
\hline
Seguridad (\textit{Security}) & \texttt{risk-engine}, \texttt{ai-client}, \texttt{audit} & Motor de evaluaci\'on de riesgos, cliente de IA generativa y registro de auditor\'ia. \\
\hline
\textit{Middleware} & \textit{Middleware} transversal & Autenticaci\'on, autorizaci\'on, manejo de errores y validaci\'on de esquemas. \\
\hline
Utilidades (\textit{Utils}) & Funciones auxiliares & Formateo de datos, generaci\'on de tokens, cifrado y funciones de apoyo. \\
\hline
\end{tabular}
\end{table}

\subsection{An\'alisis de Resultados por Componente}
\label{subsec:analisis-unitarias}

La categor\'ia de \textbf{rutas} (\textit{routes}) concentra la mayor cantidad
de casos de prueba, dado que cada \textit{endpoint} de la API REST es verificado
contra m\'ultiples escenarios: solicitudes v\'alidas con respuestas exitosas,
solicitudes con par\'ametros inv\'alidos o ausentes, intentos de acceso sin
autenticaci\'on, y operaciones con permisos insuficientes. Los suites
correspondientes a los m\'odulos \texttt{iam} y \texttt{auth} presentan la
mayor densidad de casos de prueba, lo cual refleja la complejidad inherente a
las operaciones de gesti\'on de identidad y autenticaci\'on adaptativa.

Los componentes de \textbf{seguridad} (\textit{security}) incluyen pruebas
especializadas para el motor de evaluaci\'on de riesgos (\texttt{risk-engine}),
el cual implementa algoritmos de puntuaci\'on basados en m\'ultiples factores
contextuales como la geolocalizaci\'on, el dispositivo, el horario de acceso y
el historial de comportamiento del usuario. El suite \texttt{ai-client} verifica
la integraci\'on con el servicio de inteligencia artificial generativa utilizado
para la s\'intesis de m\'etricas de seguridad, incluyendo escenarios de
\textit{timeout}, respuestas malformadas y manejo de errores de la API externa.

\begin{figure}[htbp]
\centering
\includegraphics[width=0.95\textwidth]{uniteback}
\caption{Resultados de ejecuci\'on de las pruebas unitarias del servidor \textit{backend}}
\label{fig:pruebas-unitarias-backend}
\end{figure}

La Figura~\ref{fig:pruebas-unitarias-backend} presenta la salida completa de
la ejecuci\'on de Jest, donde se observa el estado de aprobaci\'on de cada
suite individual, la cantidad de casos de prueba por suite y el tiempo de
ejecuci\'on parcial. La ausencia total de fallos confirma la estabilidad del
c\'odigo base y la correcta implementaci\'on de los contratos funcionales
definidos en los requisitos del sistema.

El tiempo total de ejecuci\'on de \textbf{19.375 segundos} para 334 casos de
prueba arroja un promedio de aproximadamente 58 milisegundos por caso de prueba,
lo cual demuestra la eficiencia del enfoque de \textit{mocking} y aislamiento
de dependencias empleado en la estrategia de pruebas. Este rendimiento permite
la integraci\'on de las pruebas unitarias en el flujo de integraci\'on continua
(CI) sin impactar significativamente los tiempos de construcci\'on del proyecto.
